\chapter{Foundations}

\section{Transcript-Limited Refinement}

A refinement system is a sequence of state updates together with a sequence of observable transcript variables.  The state space records what the process may know, while the transcript records what has actually been exposed.

\begin{definition}[Refinement system]
A refinement system is a tuple
\[
(\mathcal S, X, (R_t)_{t=1}^T, (T_t)_{t=1}^T)
\]
where \(\mathcal S\) is a state space, \(X\) is the target variable, \(R_t\) is the \(t\)-th refinement rule, and \(T_t\) is the transcript emitted at step \(t\).
\end{definition}

\section{Capacity Bound}

The basic capacity assumption is that each refinement step has bounded conditional information gain:
\[
I(X;T_t \mid T_{t-1}) \le C_t .
\]
For \(t=1\), the expression is interpreted with \(T_0\) equal to the empty transcript.

\begin{lemma}[Finite transcript budget]
If \(I(X;T_t \mid T_{t-1}) \le C_t\) for all \(t\), then the total transcript gain is bounded by
\[
I(X;T_1,\ldots,T_T) \le \sum_{t=1}^T C_t .
\]
\end{lemma}

\begin{proof}
Apply the chain rule for mutual information and then apply the per-step capacity bound.
\end{proof}

\section{Status}

This chapter fixes the language of bounded refinement.  It does not by itself prove any domain-specific lower bound, physical theorem, or Clay-problem result.
